% ============================================================================
% INTRODUCTION GÉNÉRALE
% ============================================================================

L'essor spectaculaire de l'intelligence artificielle au cours de la dernière décennie a
profondément transformé la manière dont les êtres humains interagissent avec les systèmes
numériques. Depuis l'émergence des grands modèles de langage (LLM) en 2022, incarnés par des
architectures telles que GPT, LLaMA et Gemini, une nouvelle ère d'interaction homme-machine
s'est ouverte, où l'utilisateur peut exprimer ses besoins en langage naturel et obtenir des
réponses contextualisées, structurées et actionnables. Cette révolution technologique offre
des opportunités considérables pour l'industrie touristique, un secteur intrinsèquement
dépendant de la personnalisation et de la gestion de contraintes multidimensionnelles.

L'île de Djerba, inscrite au patrimoine mondial de l'UNESCO depuis 2023, constitue l'une des
destinations touristiques les plus emblématiques de la Tunisie. Avec environ 2,5~millions de
visiteurs annuels, cette île méditerranéenne offre un tissu touristique d'une richesse
exceptionnelle~: stations balnéaires bordant des plages de sable fin, héritage culturel amazigh
préservé dans les villages traditionnels, médinas historiques, et une gastronomie
méditerranéenne authentique. Cette diversité, si elle constitue un atout majeur, rend également
la navigation et la planification d'un séjour particulièrement complexes pour le visiteur
néophyte, confronté à une offre abondante mais fragmentée.

L'état actuel des outils de planification touristique révèle des lacunes structurelles
significatives. Les moteurs de recherche traditionnels retournent des listes d'établissements
classés par popularité, sans composer un itinéraire cohérent. Les plateformes de réservation
traitent chaque prestation de manière isolée~--- hébergement, activités, restauration~--- sans
orchestrer ces éléments en un plan multi-journées adapté aux contraintes spécifiques du
voyageur.

C'est dans ce contexte que s'inscrit le présent projet de fin d'études, qui propose
\textbf{Guidni}, une plateforme de planification de voyage intelligente articulée autour de
deux sous-systèmes complémentaires. Le premier, l'\textbf{Agent Planificateur}, est un agent
conversationnel construit sur une machine à états LangGraph à quatre nœuds, disposant de
dix-huit outils appelables par le LLM et d'un pipeline RAG multilingue pour la recherche
sémantique. Le second, le \textbf{Moteur de Recommandation Hybride LinUCB}, exploite un bandit
contextuel avec un vecteur de caractéristiques de dimension~306 pour recommander des contenus
touristiques en temps réel.

Le présent rapport est organisé comme suit~:

\begin{itemize}
    \item Le \textbf{chapitre~1} présente le contexte général du projet~: l'organisme d'accueil,
    la présentation du projet, l'étude de l'existant, la solution proposée et le choix
    méthodologique (Scrum).

    \item Le \textbf{chapitre~2} est consacré au lancement du projet~: l'analyse des besoins
    fonctionnels et non fonctionnels, les définitions théoriques, l'architecture fonctionnelle
    et logique, le Product Backlog, la planification des Sprints et l'environnement de travail.

    \item Le \textbf{chapitre~3} détaille le Sprint~1~: mise en place de l'environnement de
    développement, interface utilisateur et architecture de base.

    \item Le \textbf{chapitre~4} détaille le Sprint~2~: implémentation du pipeline RAG
    multilingue et recherche sémantique.

    \item Le \textbf{chapitre~5} détaille le Sprint~3~: agent LLM, machine à états LangGraph
    et planification conversationnelle.

    \item Le \textbf{chapitre~6} détaille le Sprint~4~: moteur de recommandation LinUCB, tests
    et évaluation.
\end{itemize}

Enfin, la \textbf{conclusion générale} synthétise les contributions du projet, discute les
limitations identifiées et propose des perspectives d'évolution future.
